\section{SVD 给出最佳低秩近似}
\label{sec:13-Machine-Learning/04-Dimensionality-Reduction/02}

你在压缩一个已经训练好的模型。某个全连接层的权重是 \(W\in\R^{4096\times 4096}\)，一千七百万个参数，占了整个模型体积的三分之一。有人提议把它换成两个瘦矩阵的乘积 \(W\approx BC^{\trans}\)，其中 \(B,C\in\R^{4096\times 64}\)，参数量掉到原来的三十分之一。做法是对 \(W\) 做 SVD、取前 64 项。评审会上有人问了两句：这是最好的 64 秩近似吗？如果我还要求 \(B\) 非负、或者两个因子都量化成 8 比特，截断还是最好的吗？第一问有干净的答案，第二问的答案是否定的，而两问之间的距离正是这一节要讲的东西。

\subsection{奇异值是一张作用强度清单}

上一节从样本方差进入 PCA，那里的 \(X_c\) 是一张数据表。这一节的 \(A\) 不必是数据表，它可以是任何线性映射：权重矩阵、图像、算子的离散化。SVD 说的是，任何 \(A\in\R^{m\times n}\) 都能写成

\[
A = U\Sigma V^{\trans} = \sum_{j=1}^{r}\sigma_j\,u_jv_j^{\trans},
\qquad \sigma_1\ge\sigma_2\ge\cdots\ge\sigma_r>0,
\]

其中 \(r=\rank(A)\)。每一项 \(\sigma_ju_jv_j^{\trans}\) 是一条独立的作用通道：输入落在 \(v_j\) 上的分量被放大 \(\sigma_j\) 倍，送到输出侧的 \(u_j\) 方向。存在性与几何解释见\hyperref[sec:03-Linear-Algebra/07-SVD-and-Data-Applications/01]{``任意矩阵的奇异值分解''}一节。

于是 \(\sigma_1,\dots,\sigma_r\) 就是一张按强度排好序的清单。若 \(\sigma_1\gg\sigma_2\)，那么对大多数输入，\(A\) 干的事情几乎就是``投影到 \(v_1\)、放大、打到 \(u_1\)''，其余通道的贡献淹没在数值噪声里。截断的想法由此自然浮现：只留清单的前 \(k\) 条，

\[
A_k = \sum_{j=1}^{k}\sigma_ju_jv_j^{\trans},
\]

它的秩不超过 \(k\)，而且这个式子已经是现成的两因子形式，直接对应上面那个 \(BC^{\trans}\)。

\subsection{截断用尽了秩预算，所以没有别的矩阵能做得更好}

Eckart--Young--Mirsky 定理把``最好''这件事说死了。在谱范数与 Frobenius 范数下，

\[
\min_{\rank(B)\le k}\norm{A-B}_2 = \sigma_{k+1},
\qquad
\min_{\rank(B)\le k}\norm{A-B}_F = \Bigl(\sum_{j>k}\sigma_j^2\Bigr)^{1/2},
\]

两个下界都由同一个 \(A_k\) 达到。两种范数量的东西并不一样，最优解却重合，这不是显然的事。

谱范数那一半的证明可以压成三步。任取 \(\rank(B)\le k\)，则 \(B\) 的零空间维数至少是 \(n-k\)；它与 \(v_1,\dots,v_{k+1}\) 张成的 \((k+1)\) 维子空间维数相加超过 \(n\)，两者必有非零交，取其中的单位向量 \(z\)。于是

\[
\begin{aligned}
\norm{A-B}_2^2
&\ \ge\ \norm{(A-B)z}_2^2 \ =\ \norm{Az}_2^2\\
&\ =\ \sum_{j\le k+1}\sigma_j^2\,\ip{v_j}{z}^2
\ \ge\ \sigma_{k+1}^2\sum_{j\le k+1}\ip{v_j}{z}^2
\ =\ \sigma_{k+1}^2 .
\end{aligned}
\]

第二个等号用了 \(Bz=0\)，最后一步用了 \(z\) 是单位向量且落在那 \(k+1\) 个方向张成的空间内。上界方向只需验证 \(\norm{A-A_k}_2=\sigma_{k+1}\)，两头一夹即得。直觉上这段话说的是：秩不超过 \(k\) 的矩阵最多完整保存 \(k\) 条通道，必然漏掉一条强度不低于 \(\sigma_{k+1}\) 的；截断的策略是把秩预算全部花在最强的通道上，恰好顶到下界。

条件对账：论证用到的只有两件事——范数是酉不变的（谱范数、Frobenius 范数、乃至一般的 Schatten 范数都满足），以及可行集就是 \(\set{B:\rank(B)\le k}\) 本身，没有别的限制。这两条一旦松动，结论立刻失去保护。更细的推导与几何图景见\hyperref[sec:03-Linear-Algebra/07-SVD-and-Data-Applications/03]{``低秩近似：用少数方向保留主要结构''}一节。

\subsection{最优性写在范数里，更写在``无约束''三个字里}

回到评审会上的第二问。要求因子非负，可行集从 \(\set{\rank(B)\le k}\) 缩到它与非负锥的交；要求因子稀疏，缩到它与某个稀疏结构集的交；要求 8 比特量化，缩到一个有限点集。这些集合都不是酉不变的——旋转一下就跑出去了——而上面那段证明的每一步都建立在酉不变之上。

后果比``界变松了''更直接：\(A_k\) 通常根本不在新的可行集里，它连一个候选都算不上。

\begin{center}
\begin{tikzpicture}[>=stealth,scale=1.0]
% 低秩集合（曲线，非凸）
\draw[black!45,thick] plot[domain=-3.1:3.1,samples=60] ({\x},{0.17*\x*\x});
\node[font=\scriptsize,black!55] at (-2.55,1.55) {\(\set{\rank(B)\le k}\)};
% 可行子段
\draw[very thick,black] plot[domain=1.5:3.1,samples=30] ({\x},{0.17*\x*\x});
\node[font=\scriptsize,anchor=west] at (3.2,1.63) {加了约束后的可行集};
% 目标矩阵 A
\fill[black] (-0.2,2.5) circle (2pt);
\node[font=\scriptsize] at (-0.5,2.72) {\(A\)};
% A_k
\fill[black] (0.55,0.0514) circle (2pt);
\node[font=\scriptsize] at (0.62,-0.35) {\(A_k\)};
\draw[dashed,black!70] (-0.2,2.5) -- (0.55,0.0514);
\node[font=\scriptsize,black!70] at (-0.42,1.25) {\(\sigma_{k+1}\)};
% 约束下的最优点
\fill[black] (1.5,0.3825) circle (2pt);
\node[font=\scriptsize] at (1.62,0.02) {\(B^\star\)};
\draw[dashed,black!70] (-0.2,2.5) -- (1.5,0.3825);
\node[font=\scriptsize,black!70] at (1.15,1.7) {更大};
\end{tikzpicture}
\end{center}

\emph{截断落在低秩集合上的最近点；一旦可行集缩成粗线那一段，最近点换了位置，截断连候选都不是。}

图里那个 \(B^\star\) 才是加约束之后的最优解，而它一般无法由 \(A\) 的奇异值分解直接读出。非负矩阵分解要靠交替更新，稀疏字典学习要靠迭代阈值，量化因子分解往往落到组合优化里，这些问题多数是非凸的，算法给的是局部解，没有 Eckart--Young 那样的全局证书。工程上常见的做法是拿截断结果当初值，然后在约束集上继续优化——这是个好起点，但它只是起点，不是答案。

同样的道理适用于范数的选择。定理保证的是酉不变范数下的最优，如果你真正在意的是加权误差（某些行更重要）、是逐元素的最大偏差、或者是下游任务的准确率，那么最优解会转移到别的矩阵上，而且转移的幅度可能不小。数学给出了``最优''的精确含义；挑哪个含义，是建模者的判断。

\subsection{奇异值的尾巴不等于噪声}

把一张灰度图当成矩阵做截断，\(k=1\) 时只剩一道整体明暗渐变，\(k=5\) 时大块结构浮出来，\(k=20\) 时轮廓和纹理基本可辨，\(k=100\) 时肉眼已难以分辨。这个演示很有说服力，也很容易让人得出一个过头的结论：尾部的奇异值都是噪声。

反例在医学影像里随手可得。一张胸片上的小结节可能只占不到千分之一的像素，它贡献的能量在大尺度背景面前微不足道，对应的奇异值排在很靠后的位置。按能量占比截断，结节和噪声一起被扔掉，重建图看上去更干净，诊断信息却没了。Frobenius 范数保证的是平方误差总和最小，它从不承诺任何一个局部区域被保住。关心局部的任务，就不该用全局平方和来挑 \(k\)。

\subsection{谱隙决定哪些方向值得解释}

上一节在碎石图上找的那个拐点，在这里对应奇异值之间的间隙。\(\sigma_k\) 与 \(\sigma_{k+1}\) 拉开明显距离时，前 \(k\) 个奇异向量张成的子空间对小扰动稳定，换一批样本、加一点噪声，子空间不会大幅旋转。两者接近时，轻微扰动就足以让对应的奇异向量转到别处去。

极端情形是奇异值重复：此时单个奇异向量本来就不唯一，那个子空间里任何一组正交基都是合法的分解结果。所以纪律是：奇异值挨得近的时候不要解释单个方向，只解释它们共同张成的子空间。把某一列 loading 的正负号或具体数值讲成业务故事，是低秩分析里出错频率最高的一种做法。

\subsection{矩阵太大时，两类误差要分开记账}

\(4096\times 4096\) 还能直接做完整 SVD，再大就不行了。只需要前少量成分时，实际选择是子空间迭代、Lanczos，或者随机化 SVD。随机化的想法很简洁：取一个窄的随机矩阵 \(\Omega\)，算 \(Y=A\Omega\)，\(Y\) 的列空间以高概率覆盖了 \(A\) 的主要列空间；在 \(Y\) 的正交基上把 \(A\) 压成一个小矩阵，对小矩阵做精确 SVD，成本大幅下降。

要付出的是结果带随机性：换个种子，奇异值和奇异向量会略有不同。使用时要把误差拆成两笔账。一笔是算法的近似误差，由 oversampling 列数（通常比 \(k\) 多 5 到 10 列）和幂迭代次数控制，它是可以花计算量买回来的；另一笔是模型的截断误差 \(\sum_{j>k}\sigma_j^2\)，它由你选的 \(k\) 决定，再多算也消不掉。把两者混在一起看，就会出现``加大迭代次数却不见改善''的困惑——那说明限制你的是第二笔账。顺带一提，别用 \(A^{\trans}A\) 的特征分解绕道求奇异向量：显式形成 Gram 矩阵会把条件数平方，\(10^6\) 变成 \(10^{12}\)，双精度的有效位数被吞掉近一半，小奇异值方向首当其冲。

到这里，\(A\) 一直被当作完整已知的对象。真实系统里更常见的情形是矩阵有大片空洞：用户没看过的电影、传感器掉线的时段。观测不全时，``截断谁''这个问题连提都提不出来——下一节从缺失条目开始重建。
